\subsection{The limits of threshold adaptation}\label{sec:ceiling}
\subsubsection{Why the criticality rule missed its target}
Review~2 measured rule C in the detection pipeline against a 30\,\% false-positive-reduction target and found at best 10.8\,\%: an ICU-heavy segment traded 23.4\,\% more false positives for 48.6\,\% fewer missed attacks, a periphery-heavy segment the reverse, and the total error count (FP+FN between 294 and 311) barely moved. We now test the explanation offered then, that a threshold rule moves the operating point along a fixed trade-off curve and cannot move the curve.

\subsubsection{The trade-off curve}
%%E5TEXT%%

\subsubsection{A two-stage rule}
Stage 1 is rule C. Stage 2 is a flow classifier: a violation counts toward quarantine only if the classifier also flags the flow. Because no neural network runs inside ns-3, stage 2 is represented by the measured operating point of a trained branch on IoTID20 test windows (Table~\ref{tab:ensemble} pipeline, Attack positive): BiLSTM recall 0.9916 at a false-positive rate of 0.5103, transformer 0.9735 at 0.2194, full ensemble 0.9831 at 0.3529. Classifier errors are drawn either independently per flow, or once per device, which bounds the case where a classifier's mistakes stick to a device.

%%E7TABLE%%

Fig.~\ref{fig:tradeoff} shows the gated rule on the mimicry scenario. At the same 3.0\,\% of attackers missed, C alone quarantines 49\,\% of wearables, C gated by the BiLSTM 37\,\% (24.5\,\% fewer) and C gated by the transformer 26\,\% (47\,\% fewer): the gate moves the curve toward the origin, and only the transformer gate moves it by more than 30\,\% at constant missed attacks. Table~\ref{tab:twostage} gives the same rules in the baseline scenario over 20 seeds. With independent errors, gating C with the transformer cuts false-positive alerts by 33.4\,\% and quarantined wearables by 51.8\,\% relative to C alone while still blocking all 60 bots in every seed, meeting the 30\,\% target relative to C and reaching 66.9\,\% relative to the flat rule. Gating with the BiLSTM, the branch the roadmap proposed, reaches only 10.2\,\% and 19.8\,\%, because that branch flags 51\,\% of normal windows. When errors stick to devices the reductions grow (up to 78\,\%) but between 0.5 and 1.7 of the 60 bots escape in an average seed, because a bot the classifier misjudges once is misjudged on every packet. Within a rule, the threshold only chooses a point on the curve; a second, independent source of evidence is what moves it, and how much it helps depends on the false-positive rate of that evidence on normal traffic, not on its accuracy, which is above 0.92 for all three branches.

\subsubsection{Revised conclusion on rule C}
We therefore restate the Review~2 result in two parts. In the detection pipeline, the criticality weights act as per-class thresholds on a single classifier score; they reallocated errors between false positives and missed attacks and did not reduce false positives by 30\,\% at constant missed attacks, and no threshold on that one score can, because it only selects a point on that score's curve. In the network, device class is information the flat rule does not have, and there it changes the curve: at 3\,\% missed attackers C quarantines about half the wearables where the flat rule offers no operating point at all between 0 and 82\,\%, and at its default settings it never quarantined a life-critical or high-criticality device in any scenario, cost or two-stage variant; raising $\tau_0$ to 1,000\,ms, flattening the weights, or the mimicry attack at $\tau_0 \ge 600$\,ms does reach those classes. Threshold adaptation has a ceiling set by the information the threshold acts on. Lowering it further takes a second discriminator, and the two-stage result shows that what matters is that discriminator's false-positive rate on normal traffic.
